% !TeX root = ../../Book2.tex
% 纲要标题：### 1.4 从矩阵到不变量
% 纲要位置：计划纲要.md，第 849 行。

\section{从矩阵到不变量}
\label{b2:sec:0104}

一个算子可以有许多矩阵表示，分类时应先确定哪些变化来自坐标选择。本节写清相似关系，并用小型矩阵说明常见不变量各能辨认什么。最后陈述第四章将要证明的标准形结论，说明为什么需要把算子改写成多项式环的作用。

% 纲要内容（仅作写作计划，尚非教材正文）：
%
% * 换基与相似
% * 有理标准形和 Jordan 标准形的目标
% * 结构定理的证明归于第4章，避免重复证明
%
% ---
%

\subsection{换基与相似}
\label{b2:subsec:0104:01}

设 \(V\) 有限维，有序基为 \(\mathcal B=(e_1,\ldots,e_n)\)。向量 \(v\) 的坐标写为列向量 \([v]_{\mathcal B}\)。若算子 \(T\) 的矩阵为 \(A\)，其第 \(j\) 列就是 \([Te_j]_{\mathcal B}\)，因而
\[
[Tv]_{\mathcal B}=A[v]_{\mathcal B}.
\]
令另一组基 \(\mathcal B'=(e'_1,\ldots,e'_n)\) 满足 \(e'_j=\sum_iP_{ij}e_i\)。由于两组都是基，\(P\) 可逆，且 \([v]_{\mathcal B}=P[v]_{\mathcal B'}\)。于是
\[
[Tv]_{\mathcal B'}=P^{-1}AP[v]_{\mathcal B'}.
\]

\begin{definition}\label{b2:def:matrix-similarity}
设 \(k\) 是域，\(A,B\in M_n(k)\)。若存在 \(P\in\operatorname{GL}_n(k)\) 使 \(B=P^{-1}AP\)，就称 \(A,B\) \term[xiangsi]{相似}[similar]。同一算子在不同有序基下的矩阵恰好构成一个相似类。
\end{definition}

这里左右两侧的变换由同一个换基矩阵联系起来。对于一般线性映射 \(V\to W\)，定义域和陪域可以分别换基。若 \(A,B\in M_{m\times n}(k)\) 满足 \(B=Q^{-1}AP\)，其中 \(P\in\operatorname{GL}_n(k)\)、\(Q\in\operatorname{GL}_m(k)\)，就称 \(A,B\) 矩阵等价。它与算子的相似分类不同：任意初等行、列操作虽然保持矩阵的秩，却未必保持它代表的算子的相似类。

\begin{proposition}\label{b2:prop:similarity-invariants}
设 \(A,B\in M_n(k)\) 是域 \(k\) 上的相似矩阵。它们具有相同的秩、最小多项式、特征多项式及各个 \(p\in k[t]\) 所给的核维数 \(\dim_k\ker p(A)\)。这里把矩阵 \(A\) 的\term[tezheng-duoxiangshi]{特征多项式}[characteristic polynomial]定义为 \(\chi_A(t)=\det(tI_n-A)\)。
\end{proposition}
\begin{proof}
若 \(B=P^{-1}AP\)，归纳得 \(B^j=P^{-1}A^jP\)，故 \(p(B)=P^{-1}p(A)P\)。映射 \(v\mapsto Pv\) 将 \(\ker p(B)\) 同构到 \(\ker p(A)\)，并将 \(\operatorname{im}B\) 同构到 \(\operatorname{im}A\)，所以相应维数不变。又有 \(p(B)=0\iff p(A)=0\)，最小多项式的唯一性给出相等。

在交换环 \(k[t]\) 上使用行列式的乘法公式，得到
\[
\det(tI_n-B)=\det\bigl(P^{-1}(tI_n-A)P\bigr)=\det(tI_n-A).
\]
因此特征多项式也不变。
\end{proof}

这里使用读者已学过的行列式展开及乘法公式；行列式的外代数解释将在\chapref{b2:ch:09}建立。\(\chi_T\) 因换基不变而也可作为算子本身的记号。但单独知道特征多项式通常不足以确定相似类。例如
\[
A=\begin{pmatrix}0&0\\0&0\end{pmatrix},\qquad
B=\begin{pmatrix}0&1\\0&0\end{pmatrix}
\]
都有特征多项式 \(t^2\)，而 \(m_A=t,m_B=t^2\)，故不相似。秩也相同的例子可以在四维中得到：两个大小为二的零特征值 Jordan 块，与一个大小为三及一个大小为一的块，秩均为二、特征多项式均为 \(t^4\)，最小多项式却分别为 \(t^2,t^3\)。下面会明确这些块的形式。

\subsection{有理标准形与 Jordan 标准形}
\label{b2:subsec:0104:02}

先考察循环向量能提供什么。给定次数 \(d\ge1\) 的首一多项式
\[
p(t)=t^d+a_{d-1}t^{d-1}+\cdots+a_1t+a_0,
\]
其\term[you-juzhen]{友矩阵}[companion matrix]定义为
\[
C(p)=
\begin{pmatrix}
0&0&\cdots&0&-a_0\\
1&0&\cdots&0&-a_1\\
0&1&\cdots&0&-a_2\\
\vdots&\vdots&\ddots&\vdots&\vdots\\
0&0&\cdots&1&-a_{d-1}
\end{pmatrix};
\]
逐列说，\(C(p)e_j=e_{j+1}\ (j<d)\)，而 \(C(p)e_d=-\sum_{i=0}^{d-1}a_ie_{i+1}\)；这也包含 \(d=1\) 时的矩阵 \((-a_0)\)。友矩阵与由代数余子式矩阵转置得到的\term[ban-sui-juzhen]{伴随矩阵}[adjugate matrix]是不同对象；本书分别使用这两个名称。

\ProseParagraphBreak
若 \(v\) 为循环向量，且 \(m_T=p\)，则\cref{b2:prop:cyclic-basis}给出基 \(v,Tv,\ldots,T^{d-1}v\)。前 \(d-1\) 个基向量各被送到后一个，而
\[
T^dv=-a_0v-a_1Tv-\cdots-a_{d-1}T^{d-1}v,
\]
所以算子在这组基下的矩阵就是 \(C(p)\)。这已完整解释单个循环子空间的矩阵形式。一般算子怎样分成循环部分，以及不同分法为什么给出相同的分类数据，仍需要证明。

\begin{claim}[有理标准形，证明见第四章]\label{b2:claim:rational-canonical-goal}
设 \(V\ne0\) 是域 \(k\) 上的有限维向量空间，\(T\in\operatorname{End}_k(V)\)。在适当的有序基下，\(T\) 的矩阵为
\[
\operatorname{diag}\bigl(C(d_1),\ldots,C(d_r)\bigr),
\qquad d_1\mid d_2\mid\cdots\mid d_r,
\]
其中 \(d_i\in k[t]\) 首一且次数为正，\(\sum_i\deg d_i=\dim_k V\)。这些多项式唯一，且 \(m_T=d_r\)、\(\chi_T=\prod_i d_i\)。
\end{claim}

这一形式称为\term[youli-biaozhunxing]{有理标准形}[rational canonical form]。“有理”不要求 \(k=\mathbb Q\)，而是说所有构造都在原基域上进行，不要求先把多项式分解成一次因子。存在性与唯一性将在\cref{b2:thm:rational-canonical-form}证明，最小多项式与特征多项式公式见\cref{b2:prop:operator-characteristic-minimal}；本章的证明与习题均不依赖这一前瞻陈述。

\ProseParagraphBreak
若多项式在 \(k\) 上分裂，还可使用更细的块。定义
\[
J_d(\lambda)=
\begin{pmatrix}
\lambda&1&&0\\
&\lambda&\ddots&\\
&&\ddots&1\\
0&&&\lambda
\end{pmatrix},\qquad J_1(\lambda)=(\lambda).
\]
它称为\term[Jordan-kuai]{Jordan 块}[Jordan block]。减去 \(\lambda I\) 后，每个 \(e_j\ (j>1)\) 被送到 \(e_{j-1}\)，\(e_1\) 被送到零；因而第 \(r\) 次幂的核维数为 \(\min(r,d)\)。

\begin{claim}[Jordan 标准形，证明见第四章]\label{b2:claim:jordan-canonical-goal}
设 \(V\) 是域 \(k\) 上的有限维向量空间，\(T\in\operatorname{End}_k(V)\)。若 \(T\) 的特征多项式在 \(k\) 上分裂为一次因子的乘积，则在适当的有序基下，\(T\) 的矩阵是若干 Jordan 块 \(J_d(\lambda)\)（\(d\ge1,\lambda\in k\)）的直和。这些块除排列次序外唯一。
\end{claim}

这称为\term[Jordan-biaozhunxing]{Jordan 标准形}[Jordan canonical form]。完整证明归于\cref{b2:thm:jordan-canonical-form}。核维数序列说明了块大小为什么可望成为可恢复的不变量：对固定特征值，一个大小为 \(d\) 的块在第 \(r\) 次幂的核中贡献 \(\min(r,d)\) 维；相邻差便计算大小至少为 \(r\) 的块数。这是分析已给定块矩阵的直接计算，还不是每个算子都有这种分解的证明。

\subsection{标准形的模论解释}
\label{b2:subsec:0104:03}

对于单个循环子空间，多项式与向量之间有一个明确的对应。

\begin{proposition}\label{b2:prop:cyclic-polynomial-quotient}
设 \(V\ne0\) 是域 \(k\) 上的有限维向量空间，\(T\in\operatorname{End}_k(V)\)，\(0\ne v\in V\)，且 \(p=m_{T,v}\) 是零化 \(v\) 的首一最小多项式。映射
\[
k[t]/(p)\longrightarrow k[T]v,\qquad f+(p)\longmapsto f(T)v
\]
是 \(k\) 线性同构，并把左侧乘以 \(t\) 的运算变为右侧施加 \(T\) 的运算。
\end{proposition}
\begin{proof}
映射 \(k[t]\to k[T]v\)、\(f\mapsto f(T)v\) 线性且满，其核恰好是 \(p\) 的全部倍数，即理想 \((p)\)。\cref{b2:prop:linear-first-isomorphism}给出商空间到像的线性同构。对于任何 \(f\)，\((tf)(T)v=T(f(T)v)\)，证明最后的相容性。
\end{proof}

商环 \(k[t]/(p)\) 在这里同时有向量空间结构和多项式作用。带余除法说明 \(1,t,\ldots,t^{d-1}\) 的陪集是基；用这组基表示乘 \(t\)，又得到 \(C(p)\)。后续的模结构定理要做的，就是把一般有限维算子分解为这样的循环商，并识别每一部分的关系多项式。本节只证明了单个循环部分的对应，不以标准形的存在性作为它的依据。

\ProseParagraphBreak
基域对分类形式的影响可以从一个二维例子看出。在 \(\mathbb R^2\) 上取
\[
A=\begin{pmatrix}0&-1\\1&0\end{pmatrix}=C(t^2+1).
\]
因为 \(A^2=-I\)，且 \(e_1,Ae_1=e_2\) 是基，故最小多项式为 \(t^2+1\)。这个多项式在 \(\mathbb R\) 上没有根，算子没有实特征向量。把同一矩阵放在 \(\mathbb C^2\) 上，向量 \((1,-i)\)、\((1,i)\) 分别对应特征值 \(i,-i\)，而且二者独立，所以矩阵在复数域上相似于 \(\operatorname{diag}(i,-i)\)。矩阵条目没有变化，可用的标量与基向量却变了。讨论标准形时，必须同时写明所工作的域。
